Large language models are increasingly deployed with system prompts containing conditional (if-then) rules, yet no dedicated benchmark evaluates how reliably models follow such instructions. We introduce \benchname, a benchmark of 120 test cases spanning 8 conditional instruction types---keyword, format, content, length, negation, nested, conjunction, and disjunction---each with balanced trigger and non-trigger inputs and deterministic verification. We evaluate \gptlarge and \gptsmall and find that conditional instruction-following capacity varies dramatically by instruction type: \gptlarge ranges from 100\% accuracy on keyword conditions to 75\% on nested conditionals, while \gptsmall ranges from 100\% to 50\%. Critically, the difficulty ranking across instruction types is highly consistent between models (Spearman $\rho = 0.83$, $p = 0.011$), suggesting that certain conditional instruction types are inherently harder rather than reflecting model-specific weaknesses. Error analysis reveals an asymmetry in failure modes: the stronger model predominantly over-triggers (false positives), while the weaker model predominantly misses triggers (false negatives). We establish a difficulty hierarchy---keyword $<$ negation $<$ disjunction $<$ format $<$ conjunction $<$ content $<$ length $<$ nested---and show that the capacity bottleneck lies in the complexity of individual conditions, not the number of simultaneous conditions. Our findings provide actionable guidance for practitioners designing conditional system prompts.
